% interaction/bcs.tex — BCS transformation for hybrid oracle reductions

\section{The BCS Transformation}\label{sec:interaction-bcs}

The BCS (Ben-Sasson--Chiesa--Spooner) transformation converts an interactive
oracle reduction into a non-interactive argument by replacing oracle messages
with commitments.
In the IOP literature, BCS is typically stated for \emph{oracle proofs}
(output = accept/reject, no output oracles).
Here we develop BCS for \emph{oracle reductions}, where the output includes
oracle statements that downstream reductions can query.
This generalization is necessary for modular composition of SNARKs: each
sub-reduction in a composed pipeline may produce oracle outputs that the next
sub-reduction queries.

\subsection{Hybrid oracle reductions}

Recall from Section~\ref{sec:interaction-oracle} that an
$\OracleReduction$ pairs an oracle prover with a verifier whose oracle
access grows along the transcript path.
In practice, not every sender message is an oracle: some prover messages are
plain metadata (trace length, layout information, binding order for
sumcheck, etc.) that may legitimately shape the protocol tree.

\begin{definition}[HybridDecoration]
  \label{bcs:hybrid-decoration}
  A \emph{hybrid decoration} assigns an optional $\OracleInterface$ at
  each sender node:
  \[
    \mathsf{HybridDecoration} \;\defeq\;
    \mathsf{Role.Refine}\;(\lambda X.\; \Option\;\OracleInterface\;X).
  \]
  Sender nodes marked $\mathsf{some}\;\mathit{oi}$ are \emph{oracle senders}
  (queryable, candidates for commitment).
  Sender nodes marked $\mathsf{none}$ are \emph{plain senders}
  (sent in the clear, may shape the tree).
  \lean{Interaction.HybridDecoration}
  \uses{int:oracle-decoration}
\end{definition}

\begin{definition}[HybridOracleReduction]
  \label{bcs:hybrid-oracle-reduction}
  A \emph{hybrid oracle reduction} is the analog of $\OracleReduction$ using
  $\mathsf{HybridDecoration}$ instead of $\OracleDeco$.
  It consists of:
  \begin{itemize}
    \item A prover $P$ (an $\OracleProver$).
    \item A verifier $V$: a $\Counterpart.\mathsf{withMonads}$ with
      the monad decoration from
      $\mathsf{HybridDecoration.toMonadDecoration}$, which accumulates
      oracle access only at $\mathsf{some}\;\mathit{oi}$ sender nodes.
    \item An output oracle simulation $\mathit{sim}$: given a transcript
      $\mathit{tr}$, maps queries to output oracle family
      $\OStmtOut(\mathit{tr})$ into computations in
      $\OracleComp([\OStmtIn]_o + \mathsf{toOracleSpec}(\mathit{hd},
      \mathit{tr}))$.
  \end{itemize}
  \lean{Interaction.HybridDecoration.HybridOracleReduction}
  \uses{bcs:hybrid-decoration, int:oracle-reduction}
\end{definition}

\subsection{HybridSpec and commitment decoration}

To formalize BCS computably, we use $\mathsf{HybridSpec}$, a variant of
$\Spec$ with two kinds of nodes:

\begin{definition}[HybridSpec]
  \label{bcs:hybrid-spec}
  \begin{align*}
    \mathsf{HybridSpec} \;::=\;
    &\;\mathsf{done} \\
    \mid\;&\;\mathsf{branch}\;X\;(\mathit{rest} : X \to \mathsf{HybridSpec})
      & \text{(continuation depends on } x : X \text{)} \\
    \mid\;&\;\mathsf{pass}\;X\;\mathit{rest}
      & \text{(continuation is structurally constant)}
  \end{align*}
  $\mathsf{branch}$ nodes are used for plain senders (metadata) and receivers
  (challenges), where the continuation may depend on the message.
  $\mathsf{pass}$ nodes are used for oracle senders, where the continuation
  \emph{must not} depend on the message value (since BCS will hide the
  message behind a commitment).
  \lean{Interaction.HybridSpec}
  \uses{}
\end{definition}

The key property: at a $\mathsf{pass}\;X\;\mathit{rest}$ node,
$\Transcript\;\mathit{rest}.\mathsf{toSpec}$ does not depend on
$x : X$ \emph{definitionally}.
This eliminates the need for $\mathsf{Classical.arbitrary}$ or propositional
casts when projecting transcripts.

\begin{definition}[CommitDeco]
  \label{bcs:commit-deco}
  A \emph{commitment decoration} selects, at each $\mathsf{pass}$ node,
  whether to commit the oracle message ($\mathsf{some}\;\mathit{nc}$) or
  leave it in the clear ($\mathsf{none}$).
  At $\mathsf{branch}$ nodes, the selection is indexed by the message value
  (since the subtree depends on it).
  \lean{Interaction.HybridSpec.CommitDeco}
  \uses{bcs:hybrid-spec}
\end{definition}

\begin{definition}[SharedTranscript]
  \label{bcs:shared-transcript}
  The \emph{shared transcript} relative to a commitment decoration retains
  all $\mathsf{branch}$ messages and non-committed $\mathsf{pass}$ messages,
  but drops committed $\mathsf{pass}$ messages.
  This is the data visible to both the original and BCS-transformed
  protocols.
  \lean{Interaction.HybridSpec.SharedTranscript}
  \uses{bcs:commit-deco}
\end{definition}

\subsection{The BCS-transformed protocol spec}

\begin{definition}[bcsSpec]
  \label{bcs:bcs-spec}
  The BCS-transformed spec replaces each committed $\mathsf{pass}\;X$ node
  with $\mathsf{pass}\;\mathit{nc}.\mathsf{CommType}$ (the message type becomes
  the commitment type).
  Non-committed $\mathsf{pass}$ nodes and all $\mathsf{branch}$ nodes are
  unchanged.
  \lean{Interaction.HybridSpec.bcsSpec}
  \uses{bcs:commit-deco, bcs:hybrid-spec}
\end{definition}

\begin{definition}[bcsHybridDeco]
  \label{bcs:bcs-hybrid-deco}
  The \emph{BCS hybrid decoration} on $\mathsf{bcsSpec}$ reflects the
  restricted oracle access after commitment:
  \begin{itemize}
    \item Committed $\mathsf{pass}$ nodes $\to$ $\mathsf{none}$
      (commitment type has no oracle interface).
    \item Non-committed $\mathsf{pass}$ nodes $\to$
      $\mathsf{some}\;\mathit{oi}$ (retain oracle interface).
    \item $\mathsf{branch}$ sender nodes $\to$ $\mathsf{none}$
      (plain messages).
    \item $\mathsf{branch}$ receiver nodes $\to$ recurse.
  \end{itemize}
  \lean{Interaction.HybridSpec.bcsHybridDeco}
  \uses{bcs:bcs-spec, bcs:hybrid-decoration}
\end{definition}

\subsection{The public-query verifier decomposition}

The BCS verifier is decomposed into three components that together express
the \emph{public-query property}: the verifier's queries to committed oracles
depend only on publicly visible data.

\begin{definition}[PublicQueryVerifier]
  \label{bcs:public-query-verifier}
  A \emph{public-query verifier} consists of:
  \begin{enumerate}
    \item \textbf{Challenger} (Phase 1): a $\Counterpart.\mathsf{withMonads}$
      on $\mathsf{bcsSpec}$ with oracle access restricted to non-committed
      oracles via $\mathsf{bcsHybridDeco}$.
      Parametric in an accumulated oracle spec $\mathit{accSpec}$
      for composability.
    \item \textbf{Query function} (Phase 2a): a deterministic function
      $\mathit{queryFn} : \StmtIn \to
      \mathsf{SharedTranscript} \to \mathsf{OracleQueryDeco}$
      producing queries to committed oracles.
      The public-query property is implicit in the type: queries depend
      only on the shared transcript (publicly visible data).
    \item \textbf{Decision function} (Phase 2b): given the shared transcript
      and query responses, produces the output.
      Runs in $\OracleComp$ with access to external oracles, input oracle
      statements, and non-committed message oracles.
  \end{enumerate}
  \lean{Interaction.HybridSpec.PublicQueryVerifier}
  \uses{bcs:bcs-hybrid-deco, bcs:shared-transcript}
\end{definition}

\subsection{The central difficulty: output oracle simulation}
\label{ssec:bcs-simulate}

For oracle \emph{proofs} (output = accept/reject), the BCS transformation
is straightforward: the prover commits, the verifier challenges, openings
are verified, and the verifier decides.
There are no output oracles, so $\mathit{sim}$ is trivial.

For oracle \emph{reductions}, the situation is fundamentally more subtle.
The original reduction produces output oracle statements $\OStmtOut$
together with a simulation function
\begin{equation}\label{eq:sim-orig}
  \mathit{sim}_{\mathrm{orig}} :
  \forall\,\mathit{tr}.\;
  \mathsf{QueryImpl}\;[\OStmtOut(\mathit{tr})]_o\;
  \bigl(\OracleComp([\OStmtIn]_o +
    \mathsf{toOracleSpec}(\mathit{hd}, \mathit{tr}))\bigr)
\end{equation}
that answers output oracle queries using input oracles and the
protocol's own oracle messages.

After BCS, committed oracle messages are replaced by commitments.
The oracle spec $\mathsf{toOracleSpec}(\mathsf{bcsHybridDeco},
\mathit{bcsTr})$ includes only non-committed oracle messages.
If $\mathit{sim}_{\mathrm{orig}}$ queries a committed oracle to answer an
output oracle query, the BCS-transformed simulation
$\mathit{sim}_{\mathrm{bcs}}$ cannot do the same: it has lost access to
exactly the oracles it may need.

\begin{remark}[Why this difficulty does not arise for proofs]
  For an oracle proof, $\OStmtOut$ is empty (the output is a Boolean).
  Hence $\mathit{sim}$ is vacuous and the problem disappears.
  This is why the standard BCS literature (which treats only proofs)
  never encounters this issue.
\end{remark}

A natural but incorrect impulse is to ``transform'' $\StmtOut$ and
$\OStmtOut$ to account for the commitment step (e.g., replacing oracle
response types with commitment types or enriching the output statement).
This does not work generically: there is no uniform way to rewrite the
output interface in terms of commitments.
Instead, the output types are \emph{preserved exactly}: $\StmtOut$,
$\OStmtOut$, $\WitOut$ are the same for $\mathsf{BCS}(\calR)$ as for
$\calR$.
What changes is the internal mechanism by which the verifier computes
$\StmtOut$ and by which $\mathit{sim}$ answers output oracle queries.

\subsection{Resolution: composed spec with Phase 2}

The resolution is to include the Phase 2 opening protocol in the BCS
reduction's interaction spec.
The BCS reduction operates on a \emph{composed} spec:
\[
  \mathsf{bcsFullSpec}(\mathit{cd}, \mathit{opDeco})
  \;\defeq\;
  \mathsf{bcsSpec}(\mathit{cd})
  \;\mathbin{;}\;
  \mathsf{openingSpec}(\mathit{cd}, \mathit{opDeco}, -)
\]
where the semicolon denotes \emph{dependent} composition: the Phase 2 spec
depends on the Phase 1 transcript (because the query set depends on the
shared transcript via $\mathit{queryFn}$).

A full transcript of the composed spec is a pair
$(\mathit{tr}_1, \mathit{tr}_2)$ where:
\begin{itemize}
  \item $\mathit{tr}_1 : \Transcript(\mathsf{bcsSpec}(\mathit{cd}))$
    is the Phase 1 transcript (commitments, challenges, non-committed
    messages).
  \item $\mathit{tr}_2 : \Transcript(\mathsf{openingSpec}(\ldots,
    \mathit{tr}_1))$ is the Phase 2 transcript (opening proofs).
\end{itemize}

The Phase 2 transcript data provides query-response pairs for committed
oracles at the points determined by $\mathit{queryFn}$.
These query-response pairs fill the gap left by the absent committed oracle
access.

\subsection{BCS-transformed simulation}

The original reduction has simulation type
\begin{equation}\label{eq:sim-type}
  \mathit{sim}_{\mathrm{orig}} :
  \forall\,\mathit{tr}.\;
  \forall\,i : \iota_{so}.\;
  \OStmtOut(i).\Query \to
  \OracleComp([\OStmtIn]_o +
    \mathsf{toOracleSpec}(\mathit{hd}, \mathit{tr})) \;
  \OStmtOut(i).\Response
\end{equation}
where $\mathsf{toOracleSpec}(\mathit{hd}, \mathit{tr})$ includes all oracle
messages from the original protocol (both those that will be committed and
those that will not).

The BCS-transformed simulation $\mathit{sim}'$ has the \emph{same} outer
type signature, but operates on the composed full transcript
$\mathit{tr}_{\mathrm{full}} = (\mathit{tr}_1, \mathit{tr}_2)$
instead of the original transcript $\mathit{tr}$.
Its oracle environment is
$\OracleComp([\OStmtIn]_o +
  \mathsf{toOracleSpec}(\mathsf{bcsHybridDeco}, \mathit{tr}_1) +
  \mathit{phase2Oracles})$,
which includes input oracles, non-committed message oracles, and
query-response data from Phase 2.

Given an output oracle query $q_{\mathrm{out}}$ for index $i$,
$\mathit{sim}'$ constructs the answer as follows:
\begin{enumerate}
  \item Run $\mathit{sim}_{\mathrm{orig}}(q_{\mathrm{out}})$ as a
    computation, intercepting its oracle queries.
  \item When $\mathit{sim}_{\mathrm{orig}}$ queries a \emph{non-committed}
    oracle message, answer it directly via $\mathsf{bcsHybridDeco}$
    (these oracles are still present in $\mathit{tr}_1$).
  \item When $\mathit{sim}_{\mathrm{orig}}$ queries a \emph{committed}
    oracle at point $q$, look up the query-response pair $(q, v)$ in the
    Phase 2 data $\mathit{tr}_2$.
  \item Input oracle queries ($[\OStmtIn]_o$) pass through unchanged.
\end{enumerate}

Step~3 requires that $q$ actually appears in the Phase 2 query set.
This is not automatic: $\mathit{sim}_{\mathrm{orig}}$ might query a
committed oracle at a point that $\mathit{queryFn}$ did not include.

\subsection{Query completeness}

\begin{definition}[Query completeness]
  \label{bcs:query-completeness}
  A public-query verifier decomposition
  $(\mathit{challenger}, \mathit{queryFn}, \mathit{decide})$ together with
  a simulation function $\mathit{sim}$ is \emph{query-complete} if the
  Phase 2 query set covers all committed oracle queries that $\mathit{sim}$
  might make:
  \[
    \forall\;\mathit{tr}.\;
    \forall\;q_{\mathrm{out}} : \OStmtOut(\mathit{tr}).\mathsf{Query}.\;
    \text{every committed oracle query made by }
    \mathit{sim}(q_{\mathrm{out}})
    \text{ is in }
    \mathit{queryFn}(\mathit{st})
  \]
  where $\mathit{st} = \mathsf{projectShared}(\mathit{tr})$.
\end{definition}

In practice, query completeness is straightforward to verify: the
verifier's own queries (for $\mathit{decide}$) and the simulation queries
(for $\mathit{sim}$) are both determined from the shared transcript.
The query function $\mathit{queryFn}$ is defined to cover both.

\begin{remark}[Query completeness for proofs]
  When $\OStmtOut$ is empty (oracle proofs), query completeness is
  vacuously true.
  The query function only needs to cover the verifier's own queries, which
  is the standard public-query condition from the IOP literature.
\end{remark}

\subsection{The full BCS oracle reduction}

\begin{definition}[BCS reduction]
  \label{bcs:bcs-reduction}
  Given:
  \begin{itemize}
    \item A hybrid oracle reduction $\calR = (P, V, \mathit{sim})$ on
      $\mathit{hs}.\mathsf{toSpec}$ with hybrid decoration $\mathit{hd}$,
      input/output oracle families $\OStmtIn$, $\OStmtOut$,
      statement types $\StmtIn$, $\StmtOut$,
      and simulation $\mathit{sim}$.
    \item A commitment decoration $\mathit{cd}$ (which oracle senders to
      commit).
    \item An opening decoration $\mathit{opDeco}$ (opening protocols for
      each committed sender).
    \item A public-query verifier decomposition
      $(\mathit{challenger}, \mathit{queryFn}, \mathit{decide})$.
    \item Query completeness (Definition~\ref{bcs:query-completeness}).
  \end{itemize}
  The BCS-transformed reduction
  $\mathsf{BCS}(\calR) = (P', V', \mathit{sim}')$ is a
  $\mathsf{HybridOracleReduction}$ on
  $\mathsf{bcsFullSpec}(\mathit{cd}, \mathit{opDeco})$ with:
  \begin{itemize}
    \item \textbf{Context}: $\mathsf{bcsFullSpec}$ (composed Phase 1 + Phase 2).
    \item \textbf{Hybrid decoration}: $\mathsf{bcsHybridDeco}$ extended to
      the full composed spec.
    \item \textbf{Prover} $P'$: commits to oracle messages (Phase 1),
      then provides openings at all queried points (Phase 2).
      The prover knows the original oracle messages, so it can produce the
      same $\StmtOut$, $\OStmtOut$, and $\WitOut$ data as the original
      prover.
    \item \textbf{Verifier} $V'$: runs the challenger (Phase 1), verifies
      openings (Phase 2), then runs $\mathit{decide}$.
      The verifier computes the same $\StmtOut$ as the original, but
      using Phase 2 query-response data in place of direct committed oracle
      access.
    \item \textbf{Simulation} $\mathit{sim}'$:
      runs $\mathit{sim}_{\mathrm{orig}}$ with committed oracle queries
      answered from Phase 2 data (Subsection~\ref{ssec:bcs-simulate}).
  \end{itemize}
  Crucially, the output interface is preserved: $\StmtOut$, $\OStmtOut$,
  and $\WitOut$ are the same types as in the original reduction $\calR$.
  There is no generic way to ``transform'' the output types to account for
  commitments, and none is needed.
  What changes is the \emph{implementation}: how the verifier computes
  $\StmtOut$ and how $\mathit{sim}'$ answers output oracle queries.
  \uses{bcs:public-query-verifier, bcs:query-completeness,
    bcs:hybrid-oracle-reduction, bcs:bcs-spec}
\end{definition}

\begin{remark}[Phase 2 is protocol-internal]
  The Phase 2 opening protocol is part of the BCS reduction's interaction
  spec, not part of the output interface.
  From the perspective of the downstream consumer, the output oracle
  interface is unchanged: queries to $\OStmtOut$ are answered by
  $\mathit{sim}'$, which internally uses Phase 2 data.
  The opening mechanism is transparent to the downstream reduction.
\end{remark}

\begin{remark}[Specialization to proofs]
  For oracle proofs, $\StmtOut = \Bool$ and $\OStmtOut$ is empty.
  The output is trivially preserved (a Boolean is a Boolean).
  The verifier's $\mathit{decide}$ computes the same accept/reject decision,
  just using opening data instead of direct oracle queries.
  This matches the standard BCS formulation.
\end{remark}

\subsection{Security}

\begin{definition}[BCS completeness]
  \label{bcs:completeness}
  If the original reduction $\calR$ is complete with error
  $\varepsilon_{\mathrm{orig}}$ and the commitment schemes are correct
  (openings succeed on honestly committed data), then
  $\mathsf{BCS}(\calR)$ is complete with error
  $\varepsilon_{\mathrm{orig}} + \varepsilon_{\mathrm{commit}}$,
  where $\varepsilon_{\mathrm{commit}}$ accounts for commitment scheme
  completeness error.
  \uses{bcs:bcs-reduction, int:completeness}
\end{definition}

\begin{definition}[BCS soundness]
  \label{bcs:soundness}
  If the original reduction $\calR$ has round-by-round soundness with
  error $\varepsilon_{\mathrm{rbr}}$ and the commitment schemes satisfy
  \emph{function binding} with error $\varepsilon_{\mathrm{bind}}$ at each
  committed oracle, then $\mathsf{BCS}(\calR)$ has round-by-round
  soundness with error
  $\varepsilon_{\mathrm{rbr}} + Q \cdot \varepsilon_{\mathrm{bind}}$,
  where $Q$ is the total number of committed oracle queries (from
  $\mathit{queryFn}$).

  The function-binding property ensures that each commitment is bound to a
  specific function, so the original reduction's soundness analysis applies
  to the decommitted oracle messages.
  The error $Q \cdot \varepsilon_{\mathrm{bind}}$ accounts for the
  possibility that the adversary produces a commitment that opens
  inconsistently.
  \uses{bcs:bcs-reduction, int:rbr-soundness}
\end{definition}

\begin{remark}[Comparison with BCS for proofs]
  \label{bcs:comparison-with-proofs}
  When $\OStmtOut$ is empty (oracle proofs), the formulation above
  specializes to the standard BCS transformation:
  \begin{itemize}
    \item Query completeness is vacuous.
    \item $\mathit{sim}'$ is trivial.
    \item The composed spec reduces to $\mathsf{bcsSpec}$ plus
      the verifier's own opening queries.
    \item Soundness is the standard BCS bound from
      \cite{BCS16}.
  \end{itemize}
  The generalization to reductions adds two new elements: query
  completeness (ensuring Phase 2 covers simulation queries) and the
  composed spec (incorporating Phase 2 data into the transcript).
\end{remark}

\begin{remark}[Relationship to Chiesa--Di--Hu--Zheng]
  \label{bcs:cdh-comparison}
  The recent work of Chiesa, Di, Hu, and Zheng
  \cite{ChiesaDiHuZheng2025}
  extends BCS to interactive oracle reductions and establishes
  post-quantum security.
  Their formulation does not use $\mathsf{HybridSpec}$ or dependent
  interaction trees, but the core insight is the same: the BCS
  transformation must account for output oracle simulation, and the
  query set must be complete with respect to both the verifier's queries
  and the simulation queries.
  Our formulation makes this explicit through the query completeness
  condition and the composed spec.
\end{remark}
